\begingroup
\let\cleardoublepage\relax


\chapter*{Kontrollstrukturen in Java}
\label{ch:kontrollstrukturen-in-java}
\addcontentsline{toc}{chapter}{Kontrollstrukturen in Java}
\endgroup


\section*{If-Bedingungen}
\addcontentsline{toc}{section}{If-Bedingungen}

\subsection*{Warum braucht man If-Bedingungen?}\\
Mit If-Bedingungen kann dein Programm Entscheidungen treffen.
Zum Beispiel: "Wenn das Alter größer als 18 ist, dann zeige 'volljährig' an."

\subsubsection*{Wann brauchst du If?}
\begin{itemize}
	\item Wenn du prüfen willst, ob eine Zahl größer, kleiner oder gleich einem Wert ist
	\item Wenn du zwei Werte vergleichen willst
	\item Wenn du auf Benutzereingaben reagieren willst
	\item Wenn du Fehler vermeiden willst (z.B. prüfen, ob ein Array-Index gültig ist)
\end{itemize}

% ============================================================
\newpage


\subsection*{Einfache If-Bedingung}

\subsubsection*{Syntax:}
\begin{lstlisting}
if (Bedingung) {
    Anweisung;
}
\end{lstlisting}

\subsubsection*{Beispiele:}
\begin{lstlisting}
// Zahl pruefen
int alter = 18;
if (alter >= 18) {
    System.out.println("Volljaehrig");
}

// Array-Index pruefen
int[] zahlen = {10, 20, 30};
int index = 1;
if (index >= 0 && index < zahlen.length) {
    System.out.println(zahlen[index]);  // Sicher zugreifen
}

// Zeichen pruefen
String text = "Hallo";
if (text.charAt(0) == 'H') {
    System.out.println("Beginnt mit H");
}

// String-Laenge pruefen
String wort = "Test";
if (wort.length() > 3) {
    System.out.println("Wort ist laenger als 3 Zeichen");
}

// Wert gefunden
boolean gefunden = false;
int[] liste = {5, 10, 15, 20};
int gesucht = 15;
for (int i = 0; i < liste.length; i++) {
    if (liste[i] == gesucht) {
        gefunden = true;
    }
}
\end{lstlisting}

% ============================================================
\newpage


\subsection*{If-Else-Bedingung}

\subsubsection*{Syntax:}
\begin{lstlisting}
if (Bedingung) {
    Anweisung1;
} else {
    Anweisung2;
}
\end{lstlisting}

\subsubsection*{Beispiele:}
\begin{lstlisting}
// Gerade oder ungerade
int zahl = 7;
if (zahl % 2 == 0) {
    System.out.println("Gerade");
} else {
    System.out.println("Ungerade");
}

// Positiv oder negativ
int wert = -5;
if (wert >= 0) {
    System.out.println("Positiv oder Null");
} else {
    System.out.println("Negativ");
}

// Array leer oder gefuellt
int[] zahlen = {1, 2, 3};
if (zahlen.length > 0) {
    System.out.println("Array hat " + zahlen.length + " Elemente");
} else {
    System.out.println("Array ist leer");
}

// Zeichen vorhanden
String text = "Java";
if (text.length() > 0) {
    char erstes = text.charAt(0);
    System.out.println("Erstes Zeichen: " + erstes);
} else {
    System.out.println("String ist leer");
}
\end{lstlisting}

% ============================================================
\newpage


\subsection*{If-Else-If-Bedingung}

\subsubsection*{Syntax:}
\begin{lstlisting}
if (Bedingung1) {
    Anweisung1;
} else if (Bedingung2) {
    Anweisung2;
} else {
    Anweisung3;
}
\end{lstlisting}

\subsubsection*{Beispiele:}
\begin{lstlisting}
// Notenbewertung
int punkte = 85;
if (punkte >= 90) {
    System.out.println("Note: Sehr gut");
} else if (punkte >= 80) {
    System.out.println("Note: Gut");
} else if (punkte >= 70) {
    System.out.println("Note: Befriedigend");
} else {
    System.out.println("Note: Nicht bestanden");
}

// Groessenvergleich
int a = 5;
int b = 10;
if (a > b) {
    System.out.println("a ist groesser");
} else if (a < b) {
    System.out.println("b ist groesser");
} else {
    System.out.println("a und b sind gleich");
}

// String-Laenge kategorisieren
String eingabe = "Hallo";
if (eingabe.length() == 0) {
    System.out.println("Leer");
} else if (eingabe.length() < 5) {
    System.out.println("Kurz");
} else if (eingabe.length() < 10) {
    System.out.println("Mittel");
} else {
    System.out.println("Lang");
}

// Index-Position bestimmen
int[] werte = {10, 20, 30, 40, 50};
int index = 2;
if (index < 0) {
    System.out.println("Index zu klein");
} else if (index >= werte.length) {
    System.out.println("Index zu gross");
} else {
    System.out.println("Wert: " + werte[index]);
}
\end{lstlisting}

% ============================================================
\newpage


\subsection*{Verschachtelte If-Bedingungen}

\subsubsection*{Syntax:}
\begin{lstlisting}
if (Bedingung1) {
    if (Bedingung2) {
        Anweisung;
    }
}
\end{lstlisting}

\subsubsection*{Beispiele:}
\begin{lstlisting}
// Doppelte Pruefung
String text = "Test";
if (text.length() > 0) {
    if (text.charAt(0) == 'T') {
        System.out.println("Beginnt mit T und ist nicht leer");
    }
}

// Array-Zugriff mit Pruefung
int[] zahlen = {5, 10, 15, 20};
int index = 2;
if (index >= 0 && index < zahlen.length) {
    if (zahlen[index] > 10) {
        System.out.println("Wert an Position " + index + " ist groesser als 10");
    }
}

// Mehrere Bedingungen
int zahl = 24;
if (zahl > 0) {
    if (zahl % 2 == 0) {
        if (zahl < 100) {
            System.out.println("Positive gerade Zahl unter 100");
        }
    }
}
\end{lstlisting}

% ============================================================
\newpage


\section*{For-Schleifen}
\addcontentsline{toc}{section}{For-Schleifen}

\subsection*{Warum braucht man For-Schleifen?}\\
Mit For-Schleifen kannst du Code mehrmals ausführen lassen. Das ist nützlich, wenn du:
\begin{itemize}
	\item Durch alle Werte in einem Array gehen willst
	\item Jeden Buchstaben in einem Text prüfen willst
	\item Etwas eine bestimmte Anzahl wiederholen willst (z.B. 10-mal)
	\item Array-Werte ändern willst (dafür brauchst du die Position/Index)
\end{itemize}

% ============================================================
\newpage


\subsection*{Standard For-Schleife}

\subsubsection*{Syntax:}
\begin{lstlisting}
for (Initialisierung; Bedingung; Inkrement) {
    Anweisung;
}
\end{lstlisting}

\subsubsection*{Wichtig:} Du kannst alle drei Teile (Start, Bedingung, Erhöhung) auch weglassen - das ist erlaubt, aber für Anfänger meist nicht nötig.

\subsubsection*{Typisches Muster für Arrays:}
\begin{lstlisting}
for (int i = 0; i < array.length; i++) {
    // Verwende array[i]
}
\end{lstlisting}

\subsubsection*{Typisches Muster für Strings:}
\begin{lstlisting}
for (int i = 0; i < text.length(); i++) {
    // Verwende text.charAt(i)
}
\end{lstlisting}

\subsubsection*{Beispiele:}
\begin{lstlisting}
// Zahlen von 0 bis 4 ausgeben
for (int i = 0; i < 5; i++) {
    System.out.println(i);  // Gibt aus: 0, 1, 2, 3, 4
}

// Array durchlaufen
int[] zahlen = {10, 20, 30, 40, 50};
for (int i = 0; i < zahlen.length; i++) {
    System.out.println(zahlen[i]);
}

// String durchlaufen
String text = "Hallo";
for (int i = 0; i < text.length(); i++) {
    char c = text.charAt(i);
    System.out.println(c);
}

// Array-Werte aendern
int[] werte = {1, 2, 3, 4, 5};
for (int i = 0; i < werte.length; i++) {
    werte[i] = werte[i] * 2;  // Jeden Wert verdoppeln
}
// Ergebnis: {2, 4, 6, 8, 10}

// Array befuellen
int[] quadrate = new int[5];
for (int i = 0; i < quadrate.length; i++) {
    quadrate[i] = i * i;
}
// Ergebnis: {0, 1, 4, 9, 16}

// Summe berechnen
int[] liste = {5, 10, 15, 20};
int summe = 0;
for (int i = 0; i < liste.length; i++) {
    summe = summe + liste[i];
}
System.out.println("Summe: " + summe);  // Summe: 50

// Zeichen zaehlen
String wort = "Banane";
int anzahlA = 0;
for (int i = 0; i < wort.length(); i++) {
    if (wort.charAt(i) == 'a') {
        anzahlA++;
    }
}
System.out.println("Anzahl 'a': " + anzahlA);  // Anzahl 'a': 3

// Rueckwaerts durchlaufen
int[] zahlen2 = {1, 2, 3, 4, 5};
for (int i = zahlen2.length - 1; i >= 0; i--) {
    System.out.println(zahlen2[i]);  // Gibt aus: 5, 4, 3, 2, 1
}

// Nur gerade Indizes
int[] daten = {10, 20, 30, 40, 50, 60};
for (int i = 0; i < daten.length; i = i + 2) {
    System.out.println(daten[i]);  // Gibt aus: 10, 30, 50
}

// Teilstring durchsuchen (von Position 2 bis Ende)
String satz = "Hallo Welt";
for (int i = 2; i < satz.length(); i++) {
    System.out.println(satz.charAt(i));  // Gibt aus: l, l, o, , W, e, l, t
}
\end{lstlisting}

% ============================================================
\newpage


\subsection*{Optionale Teile der For-Schleife}

\subsubsection*{Alle drei Teile können weggelassen werden!}

\subsubsection*{Initialisierung weglassen}

Wenn die Variable bereits vorher deklariert wurde:

\begin{lstlisting}
// Beispiel 1: Zaehler bereits vorhanden
int i = 0;
for (; i < 5; i++) {
    System.out.println(i);
}

// Beispiel 2: Start-Index von aussen vorgegeben
int startIndex = 2;
for (; startIndex < 10; startIndex++) {
    System.out.println(startIndex);
}
\end{lstlisting}

\subsubsection*{Bedingung weglassen}

Wenn die Bedingung immer wahr sein soll (Endlosschleife oder mit \texttt{break} beenden):

\begin{lstlisting}
// Beispiel 1: Endlosschleife mit break
int[] zahlen = {10, 20, 30, 40, 50};
for (int i = 0; ; i++) {
    if (i >= zahlen.length) {
        break;  // Schleife manuell beenden
    }
    System.out.println(zahlen[i]);
}

// Beispiel 2: Suche mit vorzeitigem Abbruch
String text = "Hallo Welt";
for (int i = 0; ; i++) {
    if (i >= text.length()) {
        break;
    }
    if (text.charAt(i) == 'W') {
        System.out.println("Gefunden an Position: " + i);
        break;
    }
}
\end{lstlisting}

\subsubsection*{Inkrement weglassen}

Wenn die Erhöhung manuell im Schleifenkörper erfolgt:

\begin{lstlisting}
// Beispiel 1: Variable Schrittweite
int[] zahlen = {1, 2, 3, 4, 5, 6, 7, 8, 9, 10};
for (int i = 0; i < zahlen.length; ) {
    System.out.println(zahlen[i]);
    if (zahlen[i] % 2 == 0) {
        i = i + 1;  // Nach geraden Zahlen: 1 Schritt
    } else {
        i = i + 2;  // Nach ungeraden Zahlen: 2 Schritte
    }
}

// Beispiel 2: Bedingte Erhoehung
String text = "aabbbcccc";
for (int i = 0; i < text.length(); ) {
    char c = text.charAt(i);
    System.out.println("Zeichen: " + c);
    // Ueberspringe alle gleichen Zeichen
    while (i < text.length() && text.charAt(i) == c) {
        i++;  // Erhoehung hier statt im for-Kopf
    }
}
\end{lstlisting}

\subsubsection*{Mehrere Teile weglassen}

\begin{lstlisting}
// Alle Teile weggelassen = Endlosschleife
// for (;;) {
//     System.out.println("Laeuft unendlich");
// }

// Nur Bedingung vorhanden
int i = 0;
for (; i < 5; ) {
    System.out.println(i);
    i++;
}
\end{lstlisting}

% ============================================================
\newpage


\subsection*{Verschachtelte For-Schleifen}

\subsubsection*{Syntax:}
\begin{lstlisting}
for (int i = 0; i < n; i++) {
    for (int j = 0; j < m; j++) {
        Anweisung;
    }
}
\end{lstlisting}

\subsubsection*{Beispiele:}
\begin{lstlisting}
// Multiplikationstabelle
for (int i = 1; i <= 3; i++) {
    for (int j = 1; j <= 3; j++) {
        System.out.println(i + " * " + j + " = " + (i * j));
    }
}

// Zwei Arrays vergleichen
int[] array1 = {1, 2, 3};
int[] array2 = {2, 3, 4};
for (int i = 0; i < array1.length; i++) {
    for (int j = 0; j < array2.length; j++) {
        if (array1[i] == array2[j]) {
            System.out.println("Gemeinsamer Wert: " + array1[i]);
        }
    }
}

// String-Vergleich (naiv)
String text1 = "abc";
String text2 = "bcd";
for (int i = 0; i < text1.length(); i++) {
    for (int j = 0; j < text2.length(); j++) {
        if (text1.charAt(i) == text2.charAt(j)) {
            System.out.println("Gemeinsames Zeichen: " + text1.charAt(i));
        }
    }
}

// Pattern ausgeben
for (int i = 1; i <= 4; i++) {
    for (int j = 1; j <= i; j++) {
        System.out.print("*");
    }
    System.out.println();
}
// Ausgabe:
// *
// **
// ***
// ****
\end{lstlisting}

% ============================================================
\newpage


\section*{While-Schleifen}
\addcontentsline{toc}{section}{While-Schleifen}

\subsection*{Warum braucht man While-Schleifen?}\\
While-Schleifen benutzt du, wenn:
\begin{itemize}
	\item Du nicht weißt, wie oft du wiederholen musst
	\item Die Schleife laufen soll, bis etwas Bestimmtes passiert
	\item Du "so lange wie nötig" wiederholen willst
	\item Du nach etwas suchst und aufhören willst, sobald du es gefunden hast
\end{itemize}

\subsubsection*{Unterschied zur For-Schleife:}
\begin{itemize}
	\item For-Schleife: "Wiederhole genau 10-mal" oder "Gehe durch alle Array-Werte"
	\item While-Schleife: "Wiederhole, solange die Bedingung wahr ist"
\end{itemize}

% ============================================================
\newpage


\subsection*{While-Schleife}

\subsubsection*{Syntax:}
\begin{lstlisting}
while (Bedingung) {
    Anweisung;
}
\end{lstlisting}

\subsubsection*{Beispiele:}
\begin{lstlisting}
// Zahlen von 0 bis 4 ausgeben
int i = 0;
while (i < 5) {
    System.out.println(i);
    i++;
}

// Element in Array suchen
int[] zahlen = {10, 20, 30, 40, 50};
int gesucht = 30;
int index = 0;
boolean gefunden = false;
while (index < zahlen.length && !gefunden) {
    if (zahlen[index] == gesucht) {
        gefunden = true;
        System.out.println("Gefunden an Position: " + index);
    }
    index++;
}

// Erste Position eines Zeichens finden
String text = "Hallo Welt";
char gesuchtesZeichen = 'W';
int position = 0;
while (position < text.length()) {
    if (text.charAt(position) == gesuchtesZeichen) {
        System.out.println("Gefunden an Position: " + position);
        break;  // Schleife beenden
    }
    position++;
}

// Solange Wert unter Grenze
int summe = 0;
int zahl = 1;
while (summe < 100) {
    summe = summe + zahl;
    zahl++;
}
System.out.println("Summe: " + summe);
System.out.println("Zahlen addiert: " + (zahl - 1));

// Array durchlaufen (alternative zu for)
int[] werte = {5, 10, 15, 20};
int idx = 0;
while (idx < werte.length) {
    System.out.println(werte[idx]);
    idx++;
}

// Bedingung wird nie erfuellt -> Schleife laeuft nicht
int x = 10;
while (x < 5) {
    System.out.println("Wird nicht ausgegeben");
    x++;
}

// Countdown
int countdown = 5;
while (countdown > 0) {
    System.out.println(countdown);
    countdown--;
}
System.out.println("Start!");

// Erstes Vorkommen eines Teilstrings finden
String satz = "Der Hund spielt im Garten";
String suche = "spielt";
int pos = 0;
boolean gefundenString = false;
while (pos <= satz.length() - suche.length() && !gefundenString) {
    boolean match = true;
    for (int j = 0; j < suche.length(); j++) {
        if (satz.charAt(pos + j) != suche.charAt(j)) {
            match = false;
        }
    }
    if (match) {
        gefundenString = true;
        System.out.println("Gefunden an Position: " + pos);
    }
    pos++;
}
\end{lstlisting}

% ============================================================
\newpage


\section*{Vergleich: For vs. While}
\addcontentsline{toc}{section}{Vergleich: For vs. While}

\subsection*{Wann For-Schleife?}
\begin{lstlisting}
// Wenn die Anzahl der Durchlaeufe bekannt ist
for (int i = 0; i < 10; i++) { }

// Beim Durchlaufen von Arrays
int[] zahlen = {1, 2, 3, 4, 5};
for (int i = 0; i < zahlen.length; i++) { }

// Beim Durchlaufen von Strings
String text = "Hallo";
for (int i = 0; i < text.length(); i++) { }

// Wenn man den Index braucht
for (int i = 0; i < array.length; i++) {
    array[i] = i * 2;  // Index wird zum Aendern gebraucht
}
\end{lstlisting}

\subsection*{Wann While-Schleife?}
\begin{lstlisting}
// Wenn die Anzahl der Durchlaeufe unbekannt ist
while (summe < 1000) {
    summe = summe + wert;
}

// Bei Suche mit vorzeitigem Abbruch
while (index < array.length && !gefunden) {
    if (array[index] == gesucht) {
        gefunden = true;
    }
    index++;
}

// Bedingungsabhaengige Wiederholung
while (bedingungErfuellt()) {
    macheEtwas();
}
\end{lstlisting}

% ============================================================
\newpage


\section*{Häufige Fehler}
\addcontentsline{toc}{section}{Häufige Fehler}

\subsection*{Endlosschleifen}

\begin{lstlisting}
// FALSCH - Endlosschleife!
int i = 0;
while (i < 5) {
    System.out.println(i);
    // i++ vergessen! Schleife laeuft unendlich
}

// RICHTIG
int i = 0;
while (i < 5) {
    System.out.println(i);
    i++;  // Zaehler erhoehen
}
\end{lstlisting}

\subsection*{Falsche Bedingungen}

\begin{lstlisting}
// FALSCH - ArrayIndexOutOfBoundsException!
int[] zahlen = {1, 2, 3};
for (int i = 0; i <= zahlen.length; i++) {  // <= ist falsch!
    System.out.println(zahlen[i]);
}

// RICHTIG
for (int i = 0; i < zahlen.length; i++) {  // < ist richtig
    System.out.println(zahlen[i]);
}
\end{lstlisting}

\subsection*{Semikolon-Fehler}

\begin{lstlisting}
// FALSCH - Semikolon beendet Schleife sofort
for (int i = 0; i < 5; i++);  // Semikolon hier ist falsch!
{
    System.out.println(i);  // Wird nur einmal ausgefuehrt, i nicht bekannt
}

// RICHTIG
for (int i = 0; i < 5; i++) {
    System.out.println(i);
}
\end{lstlisting}

\subsection*{Off-by-One Fehler}

\begin{lstlisting}
String text = "Hallo";  // length = 5, Indizes: 0-4

// FALSCH
for (int i = 0; i <= text.length(); i++) {  // i = 5 ist zu viel!
    System.out.println(text.charAt(i));
}

// RICHTIG
for (int i = 0; i < text.length(); i++) {
    System.out.println(text.charAt(i));
}

// RICHTIG - Letztes Element direkt
char letztes = text.charAt(text.length() - 1);
\end{lstlisting}
